\scanpage{102}
A vrai dire, il ne serait pas difficile d'exiger une condition 5°) qui mette ensemble 3°) et 4°), quand on se donne non $X'_1\to X'_2$, mais
\diagram{\begin{tikzcd}[ampersand replacement=\&,row sep=3em,column sep=4em,
execute at end picture={
  \path (agA102.east) -- (agB102.west)
    coordinate[pos=.12] (agH102a) coordinate[pos=.35] (agH102b)
    coordinate[pos=.58] (agH102c) coordinate[pos=.82] (agH102d);
  \draw ([yshift=-.22em]agH102a) -- ([yshift=.32em]agH102b)
    -- ([yshift=-.22em]agH102b) -- ([yshift=.32em]agH102c)
    -- ([yshift=-.22em]agH102c) -- ([yshift=.32em]agH102d);
  \path (agB102.south) -- (agD102.north)
    coordinate[pos=.10] (agV102a) coordinate[pos=.32] (agV102b)
    coordinate[pos=.54] (agV102c) coordinate[pos=.76] (agV102d);
  \draw ([xshift=-.35em]agV102a) -- ([xshift=.35em]agV102a)
    -- ([xshift=-.35em]agV102b) -- ([xshift=.35em]agV102b)
    -- ([xshift=-.35em]agV102c) -- ([xshift=.35em]agV102c)
    -- ([xshift=-.35em]agV102d) -- ([xshift=.35em]agV102d);
}]
|[alias=agA102]| X'_1 \arrow[r,"\deleted{f_1}"] \arrow[dr,"f_1"] \&
|[alias=agB102]| \deleted{\unsure{Y_1}} \arrow[d,"\deleted{f}"] \\
X'_2 \arrow[r,"f_2"] \& |[alias=agD102]| Y_{\unsure{\deleted{2}|2}}
\end{tikzcd}}
avec $f_2$ \deleted{\unsure{ét.}} induisant un hom \deleted{\unsure{fibrant}} surjectif sur les $\pi_1$, de sorte que $\widehat{\Pi}_1(Y,f_2(I'_2))$ peut s'interpréter comme un quotient de $\widehat{\Pi}_1(X'_2,I'_2)$, de sorte que l'action \deleted{de} $\Theta'_2$ de $\Theta$ sur ce dernier passe à $\widehat{\Pi}_1(Y,f_2(I'_2))$ \ldots

Mais je ne suis pas sûr si j'aurai besoin de cette généralisation commune de 3° et 4° (et $=$ de 2°, 3°, 4°, à vrai dire !).

Les conditions 2° à 4° sont pour moi surtout une \uline{recette} commode pour trouver des conditions nécessaires sur les paramètres virtuels qui repèrent un $u\in\Gamma$ indéterminé, pour qu'ils soient effectifs i.e. proviennent bel et bien d'un $u$. Le point de vue \inserted{ici} n'est nullement de dégager \deleted{des} \unsure{condition|conditions} néc. et suff. qui permettra de \uline{caractériser} les $\Theta$ provenant de $\Gamma$ \ldots
\mspage{101}
